\notechapter{N-20230316}{Fixed Points, Convergent Iteration, and Compactness Methods}

\section{A fixed point does not by itself explain an iteration}

A fixed point of a map \(T:X\to X\) is a point \(x^*\) satisfying
\[
  T(x^*)=x^*.
\]
If an iteration
\[
  x_{n+1}=T(x_n)
\]
converges to \(x\) and \(T\) is continuous, then
\[
  T(x)=T\left(\lim_{n\to\infty}x_n\right)
  =\lim_{n\to\infty}T(x_n)
  =\lim_{n\to\infty}x_{n+1}
  =x.
\]
Thus every convergent orbit ends at a fixed point.  The converse is false: the existence of a fixed point says nothing by itself about whether an orbit converges to it.

For example,
\[
  T(x)=-x
\]
has the fixed point \(0\), but an orbit starting from \(x_0\ne0\) alternates between \(x_0\) and \(-x_0\).  The map
\[
  T(x)=x
\]
has every point fixed, so there is no unique target for iteration.  A useful fixed-point theorem must therefore specify which mechanism supplies existence, uniqueness, and convergence.

\section{Contractions make successive points geometrically closer}

Let \((X,d)\) be a metric space.  A map \(T:X\to X\) is a contraction when there is a constant \(0<c<1\) such that
\[
  d(Tx,Ty)\le c\,d(x,y)
\]
for all \(x,y\in X\).

Starting from \(x_0\), define
\[
  x_{n+1}=T(x_n).
\]
The contraction estimate gives
\[
  d(x_{n+1},x_n)
  \le c\,d(x_n,x_{n-1})
  \le\cdots\le c^n d(x_1,x_0).
\]
For \(m>n\), the triangle inequality then yields
\[
\begin{aligned}
  d(x_m,x_n)
  &\le\sum_{k=n}^{m-1}d(x_{k+1},x_k)\\
  &\le d(x_1,x_0)\sum_{k=n}^{m-1}c^k\\
  &\le\frac{c^n}{1-c}d(x_1,x_0).
\end{aligned}
\]
The orbit is therefore Cauchy.  This is the key analytic fact; all later conclusions depend on where that Cauchy sequence is allowed to converge.

\section{Completeness supplies the point that the estimates approach}

A metric space is complete when every Cauchy sequence converges to a point of the space.

\begin{theorem}[Banach contraction principle]
Let \((X,d)\) be a nonempty complete metric space and let \(T:X\to X\) be a contraction.  Then \(T\) has a unique fixed point \(x^*\).  For every \(x_0\in X\), the iterates \(x_{n+1}=T(x_n)\) converge to \(x^*\).
\end{theorem}

\begin{proof}
The previous estimate shows that \((x_n)\) is Cauchy, so completeness gives a limit \(x^*\in X\).  A contraction is Lipschitz and hence continuous, so
\[
  T(x^*)
  =T\left(\lim x_n\right)
  =\lim T(x_n)
  =\lim x_{n+1}
  =x^*.
\]
If \(x\) and \(y\) are fixed points, then
\[
  d(x,y)=d(Tx,Ty)\le c\,d(x,y).
\]
Since \(c<1\), this forces \(d(x,y)=0\), hence \(x=y\).
\end{proof}

Completeness cannot be omitted.  On the incomplete space \((0,1)\), the map
\[
  T(x)=\frac{x}{2}
\]
is a contraction with constant \(1/2\) and maps \((0,1)\) into itself.  Its iterates satisfy
\[
  x_n=2^{-n}x_0\longrightarrow0,
\]
but \(0\notin(0,1)\).  The contraction has no fixed point in the declared space.  The estimates have located a missing boundary point.

\section{The proof gives two practical stopping bounds}

Let \(x^*\) be the fixed point.  Passing \(m\to\infty\) in the Cauchy estimate gives the a priori bound
\[
  \boxed{
  d(x_n,x^*)
  \le\frac{c^n}{1-c}d(x_1,x_0).}
\]
It predicts in advance how many iterations are sufficient.

There is also an a posteriori bound that uses the most recent step.  Since
\[
  d(x_{n+k+1},x_{n+k})
  \le c^k d(x_{n+1},x_n),
\]
we obtain
\[
\begin{aligned}
  d(x_n,x^*)
  &\le\sum_{k=0}^{\infty}d(x_{n+k+1},x_{n+k})\\
  &\le\frac{1}{1-c}d(x_{n+1},x_n).
\end{aligned}
\]
Thus
\[
  \boxed{
  d(x_n,x^*)
  \le\frac{d(x_{n+1},x_n)}{1-c}.}
\]
This gives a rigorous stopping rule using quantities already computed by the iteration.

\section{The equation \texorpdfstring{$x=\cos x$}{x=cos x} is a complete example}

Consider
\[
  T(x)=\cos x
\]
on the interval \([0,1]\).  Since
\[
  \cos1\le\cos x\le1,
\]
the map sends \([0,1]\) into itself.  By the mean value theorem,
\[
  |\cos x-\cos y|
  \le\sin1\,|x-y|
\]
for \(x,y\in[0,1]\).  Because \(\sin1<1\), the map is a contraction with
\[
  c=\sin1\approx0.84147.
\]
The interval is complete, so there is a unique solution
\[
  x^*=\cos x^*
\]
in \([0,1]\), and every initial value in the interval converges to it.

Starting from \(x_0=1\),
\[
  x_1=\cos1,
  \qquad
  x_2=\cos(\cos1),
  \qquad
  \ldots
\]
The limit is
\[
  x^*\approx0.7390851332.
\]
If two successive iterates differ by less than \((1-c)10^{-6}\), the a posteriori estimate guarantees an absolute error below \(10^{-6}\).  The theorem therefore supplies not only existence and uniqueness but also a certified numerical algorithm.

The choice of fixed-point formulation matters.  The equation \(f(x)=0\) can often be rewritten as \(x=T(x)\) in many ways.  One formulation may be contractive while another diverges.  Fixed-point iteration is therefore partly the art of choosing a map with a small Lipschitz constant on an invariant complete set.

\section{Picard iteration moves the theorem to a function space}

Consider the initial-value problem
\[
  y'(t)=f(t,y(t)),
  \qquad
  y(t_0)=y_0.
\]
Integrating gives the equivalent equation
\[
  y(t)=y_0+\int_{t_0}^{t}f(s,y(s))\,ds.
\]
On the Banach space \(C([t_0-h,t_0+h])\) with the supremum norm, define
\[
  (Ty)(t)
  =y_0+\int_{t_0}^{t}f(s,y(s))\,ds.
\]
Assume that \(f\) is Lipschitz in its second variable:
\[
  |f(t,u)-f(t,v)|\le L|u-v|.
\]
Then
\[
\begin{aligned}
  \|Ty-Tz\|_\infty
  &\le\sup_t\int_{t_0}^{t}
    |f(s,y(s))-f(s,z(s))|\,ds\\
  &\le Lh\,\|y-z\|_\infty.
\end{aligned}
\]
If \(h<1/L\), the operator is a contraction.  One also chooses a closed ball of functions on which \(T\) maps the ball into itself; this follows from a bound on \(|f|\) on a suitable rectangle around \((t_0,y_0)\).

Banach's theorem now gives a unique fixed point of \(T\), hence a unique local solution of the differential equation.  The iterates
\[
  y_{n+1}(t)
  =y_0+\int_{t_0}^{t}f(s,y_n(s))\,ds
\]
are the Picard approximations.  The same geometric contraction estimate that solved \(x=\cos x\) now proves existence and uniqueness in an infinite-dimensional space.

\section{Topological existence does not provide Banach's algorithm}

Banach's theorem is not the only route to a fixed point.  Brouwer's theorem says that every continuous map from a nonempty compact convex subset of \(\mathbb R^n\) to itself has a fixed point.  No metric contraction is required.

The difference in conclusions is substantial.  The identity map on a closed ball satisfies Brouwer's hypotheses and has every point fixed.  A rotation of a disk fixes its center but preserves distances rather than shrinking them.  Brouwer gives existence, but not uniqueness, a convergence rate, or even a canonical iteration that approaches a fixed point.

This distinction matters in applications.  If the goal is a certified numerical method, the Banach hypotheses are valuable precisely because they are stronger.  If the map is continuous on a compact convex region but cannot be made contractive, Brouwer may still prove that a solution exists, though a separate method is needed to locate it.

\section{Schauder replaces finite-dimensional compactness by a compact map}

Closed bounded sets in infinite-dimensional Banach spaces are usually not compact, so Brouwer cannot be applied directly.  Schauder's fixed-point theorem uses compactness of the operator instead.

\begin{theorem}[Schauder fixed-point theorem]
Let \(C\) be a nonempty closed convex subset of a Banach space.  If \(T:C\to C\) is continuous and \(T(C)\) is relatively compact, then \(T\) has a fixed point.
\end{theorem}

A typical source of compactness is an integral operator.  On \(C([0,1])\), let
\[
  (Tu)(x)=g(x)+\int_0^1K(x,s)F(s,u(s))\,ds,
\]
where \(K\), \(F\), and \(g\) are continuous and a closed ball \(C\) is chosen so that \(T(C)\subseteq C\).  Uniform bounds on \(Tu\) and uniform continuity of \(K\) give equicontinuity of \(T(C)\).  Arzelà--Ascoli then shows that \(T(C)\) is relatively compact.

Schauder therefore yields a solution of the integral equation even when no Lipschitz estimate makes \(T\) contractive.  What is lost is uniqueness and convergence of the naive iteration.  Several fixed points may exist.

\section{Krasnoselskii separates the shrinking and compact parts}

Many nonlinear equations produce an operator that is neither contractive nor compact as a whole but decomposes as
\[
  T=A+B.
\]
Krasnoselskii's theorem applies on a closed bounded convex set \(C\) when \(A\) is a contraction, \(B\) is continuous and compact, and
\[
  Ax+By\in C
\]
for all \(x,y\in C\) under the relevant version of the theorem.

The decomposition reflects two different mechanisms.  The map \(A\) controls metric instability; the map \(B\) supplies compactness.  One way to see the proof strategy is to fix \(y\in C\) and solve
\[
  x=Ax+By.
\]
For fixed \(y\), the right-hand side is contractive in \(x\), so Banach gives a unique solution \(x=R(y)\).  Compactness of \(B\) is then used to show that the induced map \(R\) has a fixed point by a Schauder-type argument.

This theorem is not merely a longer list entry after Banach and Schauder.  It is designed for equations whose linear or dissipative part shrinks distances while an integral or smoothing term is compact.

When an equation has been rewritten as \(x=T(x)\), the next step is to identify which of these mechanisms is actually available.

If one can prove
\[
  d(Tx,Ty)\le c\,d(x,y),
  \qquad c<1,
\]
then Banach gives existence, uniqueness, global convergence of iteration, and explicit error bounds.  If only continuity and finite-dimensional compact convex invariance are available, Brouwer gives existence.  In an infinite-dimensional problem, relative compactness of the image points toward Schauder.  A natural decomposition into a contraction plus a compact map points toward Krasnoselskii.

The hypotheses are therefore not technical decorations around the same conclusion.  They encode different reasons a fixed point must exist, and those reasons determine what else one is entitled to claim about uniqueness, approximation, and stability.